% Part: methods
% Chapter: proofs
% Section: introduction

\documentclass[../../../include/open-logic-section]{subfiles}

\begin{document}

\olfileid[hi]{mth}{prf}{int}

\olsection{परिचय}

प्रारंभिक तर्कशास्त्र के अपने अनुभव के आधार पर आप किसी !!a{derivation}
प्रणाली---शायद प्राकृतिक निगमन या फ़िच-शैली की !!{derivation} प्रणाली, या
शायद प्रमाण-वृक्ष प्रणाली---के साथ सहज होंगे। आपको शायद याद होगा कि इन
प्रणालियों में प्रमाण करते हुए या तो कोई !!a{formula} सिद्ध करते थे या
दिखाते थे कि दिया गया तर्क वैध है। इसके लिए आप प्रणाली के नियम तब तक लागू
करते थे जब तक वांछित अंतिम परिणाम न मिल जाए। तर्कशास्त्र \emph{के बारे में}
विचार करते हुए हम भी बातें सिद्ध करते हैं, किंतु अधिकांश स्थितियों में किसी
!!a{derivation} प्रणाली का उपयोग नहीं करते। वास्तव में जिन अधिकांश प्रमाणों
पर हम विचार करते हैं वे पूर्णतः प्रथम-क्रम तर्कशास्त्र की भाषा में होने की
बजाय हिंदी में (शायद बीच-बीच में कुछ प्रतीकात्मक भाषा के साथ) किए जाते हैं।
ऐसे प्रमाण बनाते समय आरंभ में आप असमंजस में हो सकते हैं---!!a{derivation}
प्रणाली के बिना कुछ कैसे सिद्ध करूँ? आरंभ कैसे करूँ? कैसे जानूँ कि मेरा
प्रमाण सही है?

प्रमाण का प्रयास करने से पहले यह जानना महत्त्वपूर्ण है कि प्रमाण क्या है और
उसे कैसे बनाया जाए। नाम से संकेत मिलता है कि \emph{प्रमाण} यह दिखाने के लिए
होता है कि कोई बात सत्य है। इसे संवाद के रूप में सोच सकते हैं---कोई आपसे
पूछता है कि क्या कोई बात सत्य है, जैसे क्या दो के अतिरिक्त प्रत्येक अभाज्य
संख्या विषम है। केवल ``हाँ'' कहना पर्याप्त नहीं; वह जानना चाहेगा
\emph{क्यों}। ऐसी स्थिति में आप प्रमाण देंगे।

दैनिक वार्तालाप में उत्तर का केवल संकेत या अपूर्ण उत्तर पर्याप्त हो सकता
है। किंतु तर्कशास्त्र और गणित में हमें कठोर प्रमाण चाहिए---हम दिखाना चाहते
हैं कि कोई बात \emph{किसी भी} संदेह से परे सत्य है। इसका अर्थ है कि हमारे
प्रमाण का प्रत्येक चरण उचित ठहराया जाए और औचित्य ठोस हो (अर्थात् जिस मान्यता
का उपयोग कर रहे हैं वह वास्तव में सिद्ध किए जा रहे प्रमेय के कथन में मानी
गई हो, लागू की गई परिभाषाएँ सही ढंग से लागू हों, जिन औचित्यों का सहारा लिया
गया हो वे सही अनुमान हों, इत्यादि)।

सामान्यतः हम कोई कथन सिद्ध कर रहे होते हैं। सिद्ध किए जा रहे कथनों को हम
विभिन्न नाम देते हैं: प्रतिज्ञप्ति, प्रमेय, लेम्मा या उपप्रमेय। प्रतिज्ञप्ति
प्रमाण योग्य मूल कथन है: दर्ज करने योग्य महत्त्वपूर्ण, किंतु शायद विशेष गहन
नहीं अथवा प्रायः लागू नहीं। प्रमेय महत्त्वपूर्ण, सार्थक प्रतिज्ञप्ति है।
उसका प्रमाण प्रायः कई चरणों में विभाजित होता है और कभी-कभी उसका नाम उसे
पहली बार सिद्ध करने वाले व्यक्ति पर (जैसे कैंटर का प्रमेय,
लोवेनहाइम--स्कोलेम प्रमेय) अथवा उसके संबंधित तथ्य पर (जैसे पूर्णता प्रमेय)
रखा जाता है। लेम्मा ऐसी प्रतिज्ञप्ति या प्रमेय है जिसका उपयोग किसी अधिक
महत्त्वपूर्ण परिणाम के प्रमाण में होता है। भ्रमित करने वाली बात यह है कि
कभी-कभी लेम्मा अपने-आप में महत्त्वपूर्ण परिणाम होते हैं और उनका नाम भी उन्हें
प्रस्तुत करने वाले व्यक्ति पर होता है (जैसे ज़ोर्न का लेम्मा)। उपप्रमेय ऐसा
परिणाम है जो किसी अन्य परिणाम से सरलता से निष्पन्न होता है।

सिद्ध किए जाने वाले कथन में प्रायः ऐसी मान्यताएँ होती हैं जो स्पष्ट करती
हैं कि हम किस प्रकार की वस्तुओं के बारे में कुछ सिद्ध कर रहे हैं। वह
``$!A$, $!B\lif !C$ रूप का कोई !!a{formula} हो'' या ``मान लें
$\Gamma\Proves !A$'' अथवा ऐसा कुछ कहकर आरंभ हो सकता है। ये प्रतिज्ञप्ति,
प्रमेय या लेम्मा की \emph{परिकल्पनाएँ} हैं और अपने प्रमाण में आप इन्हें सत्य
मान सकते हैं। वे सिद्ध की जा रही बात को सीमित करती हैं और चर्चा की वस्तुओं
के लिए कुछ नाम भी प्रस्तुत करती हैं। उदाहरणतः यदि आपकी प्रतिज्ञप्ति
``$!A$, $!B\lif !C$ रूप का कोई !!a{formula} हो'' से आरंभ होती है, तो आप
केवल एक विशेष प्रकार के सभी सूत्रों (अर्थात् सशर्तों) के बारे में कुछ सिद्ध
कर रहे हैं और समझा जाता है कि $!B\lif !C$ वह मनमाना सशर्त है जिसके बारे में
आपका प्रमाण बात करेगा।

\end{document}
